% ============================================================
%  ch02_ostatni_dzien_wakacji.tex  —  Chapter 2
%  Ostatni dzień wakacji. Osiedle na Burowcu (pp. TBC)
% ============================================================
%
% ── STYLE RULES FOR THIS FILE ───────────────────────────────
%
%  These rules must be followed consistently throughout.
%  Drift from these rules has occurred before — check carefully.
%
%  HEADINGS:
%  - Use \section{} for ALL passage headings without exception.
%  - \subsection{} and \subsection*{} are NOT used in this file.
%  - Section titles follow the pattern:
%      Polish title / English title   (e.g. \section{Obchód / The Rounds})
%    or just a name if self-evident   (e.g. \section{Groszek})
%  - \section{} renders in blue bold with a rule beneath
%    (defined in olowiane_dzieci_analiza.tex).
%
%  POLISH SOURCE TEXT:
%  - Always use \poltext{...} — never \begin{quote}...\end{quote}
%    and never plain \textit{...} or any other wrapper.
%  - \poltext{} renders the text in red italic (keywordred).
%
%  ENGLISH TRANSLATION:
%  - Always use \poltrans{...} immediately after \poltext{}.
%  - \poltrans{} renders in blue (polishblue) with a rule beneath.
%  - No label precedes the translation — colour contrast suffices.
%
%  ANALYSIS NOTES:
%  - Follow \poltrans{} with free prose commentary.
%  - Bold Polish terms with \textbf{}, gloss in \textit{}.
%  - Standard pattern: \textbf{term} --- \textit{gloss} --- analysis.
%  - Use \medskip before thematic summary paragraphs.
%  - Grammar or vocabulary notes may be prefixed with:
%    \noindent\textit{Grammar note:} or \noindent\textit{Vocabulary note:}
%
%  PASSAGE SEPARATORS:
%  - Use % ──────────────────────────────────────────────────
%    (single % with em-dashes) between passages within a section.
%    Do NOT use the %% --- style.
%
%  UPDATE PROTOCOL:
%  - The LaTeX file is updated after each passage is confirmed
%    by the user — not in batches.
%  - Always append new content by str_replace on the final line.
%  - Verify line count after each update with wc -l.
%
% ── WHERE THE READING CURRENTLY STANDS ─────────────────────
%
%  Last passage analysed (session ending 2026-03-06):
%  Chapter 2 not yet started.
%  Chapter 1 (ch01_jeszcze_jeden_pogrzeb.tex) is complete.
%
% ============================================================


\chapter{Ostatni dzień wakacji. Osiedle na Burowcu}
\markboth{Ostatni dzień wakacji. Osiedle na Burowcu}{}

% ── Chapter title ────────────────────────────────────────────
\section{Tytuł rozdziału / Chapter Title}

\poltext{Ostatni dzień wakacji. Osiedle na Burowcu.}
\poltrans{The Last Day of the Holidays. The Estate at Burowiec.}

\textbf{Ostatni dzień wakacji} --- \textit{the last day of
the holidays}. \textbf{Ostatni} --- \textit{last, final} ---
from \textit{ostać się} (to remain); the superlative of
lateness, the day at the end of summer that everyone
who has been a schoolchild knows by feel.
\textbf{Wakacje} --- \textit{school holidays, summer
vacation} --- plural only, like \textit{drzwi} (door)
or \textit{nożyczki} (scissors): summer does not come
in the singular. The word is borrowed from Latin
\textit{vacatio} (freedom, exemption) via the academic
calendar.

\textbf{Osiedle na Burowcu} --- \textit{the estate at
Burowiec}. \textbf{Osiedle} --- the housing estate, the
\textit{bloki}: the word that has governed the spatial
world of Chapter~1. \textbf{Na Burowcu} --- locative
of \textbf{Burowiec}: a specific named estate or district,
\textit{na} indicating location on or at a place
(used with open areas, estates, squares, as distinct
from \textit{w} for enclosed spaces). Burowiec is a
district of Olkusz, a town in the Małopolska region
of southern Poland, southeast of Kraków --- the area
whose proximity to lead and zinc smelting operations
is the book's subject.

\medskip
Where Chapter~1 opened at a funeral and moved through
the present-tense reckoning of a generation's losses,
Chapter~2 announces its coordinates in time and space:
the last day of summer, a specific estate in a specific
town. The chapter will move into the past, into the
children themselves --- before the vanishing, at the
moment just before school resumes and the world of the
estate is still fully inhabited.

%% ---------------------------------------------------------------
\section{Wszystko zaczęło się od huty / It All Started with the Smelter}
%% ---------------------------------------------------------------

\poltext{Co do znikania moich rówieśników, wszystko zaczęło
się od huty. Przez sto pięćdziesiąt lat jej istnienia różnie
ją nazywano: Uthemann, Zakłady Cynkowe Szopienice, blajówa,
Huta Metali Nieżelaznych. Jak zwał, tak zwał. Tu wszystko
zaczynało i kończyło się na hucie.}

\poltrans{As for the disappearing of my contemporaries,
it all started with the smelter. Over a hundred and fifty
years of its existence it was called by various names:
Uthemann, the Szopienice Zinc Works, \textit{blajówa},
the Non-Ferrous Metals Smelter. Whatever it was called,
it was called. Here everything began and ended with the
smelter.}

The chapter opens with a declaration, a list, a proverb,
and an axiom. In four sentences Jędryk names the smelter
as the origin of everything, passes through 150 years of
its history in a list of aliases, dismisses the question
of names with a shrug, and then restates the first claim
with the finality of a closing argument. The book has
found its subject; Chapter~2 announces it without ceremony.

\textbf{Co do znikania moich rówieśników} ---
\textit{as for the disappearing of my contemporaries}.
\textbf{Co do} + genitive: a formal topic-introducer,
\textit{as for, regarding, on the matter of} ---
the register of a deposition or a report. \textbf{Znikania}
--- genitive of \textit{znikanie}, the verbal noun from
\textit{znikać}: \textit{the disappearing}, the process
nominalised. Chapter~1 used \textit{znikały} and
\textit{znikali} as verbs; now the disappearing is a
noun, a phenomenon to be accounted for. \textbf{Moich
rówieśników} --- \textit{of my contemporaries, of my
peers}: genitive plural of \textit{rówieśnik} (from
\textit{równy} --- equal, and \textit{wiek} --- age):
one who is of equal age, a peer. \textbf{Moich} ---
\textit{my} --- the possessive confirms the narrator's
place inside the cohort: not an external observer but
a member of the disappearing generation.

\textbf{wszystko zaczęło się od huty} ---
\textit{it all started with the smelter}. The thesis.
\textbf{Wszystko} --- \textit{everything, it all}.
\textbf{Zaczęło się} --- perfective reflexive past of
\textit{zaczynać się}: \textit{it began, it started}
--- the reflexive \textit{się} makes the beginning
intransitive, self-initiated, as if the thing started
of its own accord. \textbf{Od huty} --- \textit{from
the smelter}: \textit{od} + genitive of origin and
cause. The smelter is the source, the point of
departure, the \textit{od} from which all consequences
flow. The sentence could close the book as easily as
open a chapter.

\textbf{Przez sto pięćdziesiąt lat jej istnienia} ---
\textit{over a hundred and fifty years of its existence}.
\textbf{Przez} + accusative: duration, passage through
time. \textbf{Sto pięćdziesiąt lat} --- one hundred
and fifty years. \textbf{Jej istnienia} --- genitive
of \textit{jej istnienie}: \textit{its existence},
the smelter's span of operation. A century and a half:
the smelter predates the narrator's grandparents,
predates modern Poland, predates two world wars.

\textbf{różnie ją nazywano} --- \textit{it was called
by various names, people called it variously}.
\textbf{Różnie} --- \textit{variously, in different
ways}. \textbf{Ją} --- accusative of \textit{ona}:
\textit{her/it}, referring to \textit{huta} (feminine).
\textbf{Nazywano} --- impersonal passive past of
\textit{nazywać}: \textit{one called, it was called}
--- the unnamed collective of namers, the succession
of owners, regimes, and bureaucracies that renamed
it over 150 years.

\medskip
The list of names that follows is a compressed history
of the region:

\textbf{Uthemann} --- the original founder's name:
Friedrich Uthemann, a German industrialist who
established a zinc-smelting operation in Szopienice
(now a district of Katowice, in Upper Silesia) in
the nineteenth century. The German surname places the
smelter's origins in the industrial colonisation of
Silesia under Prussian rule.

\textbf{Zakłady Cynkowe Szopienice} --- \textit{the
Szopienice Zinc Works}: the interwar or early
communist-era designation, a sober administrative
title. \textbf{Zakłady} --- \textit{works, plant,
establishment} (plural of \textit{zakład}).
\textbf{Cynkowe} --- \textit{zinc} (adjectival form
of \textit{cynk}).

\textbf{blajówa} --- the workers' name, the local
colloquial name, set deliberately among the official
designations. From German \textit{Blei} (lead) +
the Polish suffix \textit{-ówa} (used to form
informal place-names and nicknames, especially
feminine): \textit{the lead-place, the leady one}.
The suffix makes it familiar, almost affectionate ---
the name the neighbourhood used, not the name on
the letterhead. Its presence in this list is the
list's most important entry: it is the only name
that tells you what was actually produced there.

\textbf{Huta Metali Nieżelaznych} --- \textit{the
Non-Ferrous Metals Smelter}: the late communist-era
official designation, bureaucratically precise and
informationally evasive. \textbf{Nieżelaznych} ---
\textit{non-ferrous}: genitive plural of
\textit{nieżelazny} (\textit{nie-} + \textit{żelazny},
iron): not iron, i.e.\ lead, zinc, cadmium. The
name describes the product category without naming
any specific metal. The more official the name, the
less it says.

\textbf{Jak zwał, tak zwał.} --- \textit{Whatever
it was called, it was called. A name's a name.}
A Polish proverb-formula: \textbf{jak\ldots{}
tak\ldots{}} sets up a parallelism of indifference.
\textbf{Zwał} is the past tense third person
singular of \textit{zwać} --- an archaic/literary
verb meaning \textit{to call, to name} --- used
here as the impersonal \textit{one called it}.
The formula means: \textit{it doesn't matter what
you called it; the thing was what it was}. The
proverb dismisses the whole question of naming
that the preceding list raised, and in doing so
makes a point: the names changed (German founder,
zinc works, workers' slang, communist designation),
but the smelter and its emissions did not.

\textbf{Tu wszystko zaczynało i kończyło się na
hucie.} --- \textit{Here everything began and ended
with the smelter.} The chapter's axiom, restating
the opening thesis with added weight. \textbf{Tu}
--- \textit{here}: the estate, the neighbourhood,
the world of these children. \textbf{Zaczynało i
kończyło się} --- imperfective reflexive past:
\textit{began and ended}, the two verbs paired and
balanced, \textit{zaczynać się} (to begin) and
\textit{kończyć się} (to end). \textbf{Na hucie}
--- \textit{at/with the smelter}: locative of
\textit{huta}, the \textit{na} of attachment and
dependence. Everything in this place had its origin
and its terminus in the smelter: employment, economy,
community --- and, the book argues, illness and death.

\medskip
\noindent\textit{Context note:} The smelter in
question --- operating under its various names in
Szopienice and the surrounding Silesian--Małopolska
region --- processed lead and zinc ore for over a
century, emitting lead particulates that settled on
the surrounding estates. The children who grew up
there in the 1970s and 1980s were breathing and
ingesting lead dust at levels now understood to
cause neurological damage, immune dysfunction, and
early death. The \textit{blajówa} --- the lead-place
--- named itself honestly. The official designations
did not.

\medskip
The opening of Chapter~2 performs the same movement
as the close of Chapter~1, but inverted: where
\textit{znikającą klasę} arrived at the end as a
discovery, \textit{wszystko zaczęło się od huty}
is stated at the start as a premise. The chapter
will now go back to show how the premise was lived.
